\subsection{Exact structural limits for the last two ordinary classes}

\begin{claim}{A.C2.24}{Proved}
Let $\alpha^2+\alpha+1=0$ and let
\[
 C_0/\mathbf F_4:\qquad
 v^2+v=z+\frac{\alpha^2(z+1)}{z^2+z+\alpha}.
                                                        \tag{C2.70}
\]
The following assertions hold for $C_0$ and its
coefficientwise-square conjugate, the two classes
left undecided by A.C2.22 and subsequently
settled by A.C2.25.
\begin{enumerate}
\item Their ordinary Jacobians are geometrically simple, with
\[
 \operatorname{End}^0_{\overline{\mathbf F}_4}(J_{C_0})
 =\mathbf Q(\pi)
 =\mathbf Q\bigl(\sqrt2,\sqrt{-7+4\sqrt2}\bigr),
                                                        \tag{C2.71}
\]
a non-Galois quartic CM field. Their geometric
automorphism group is exactly $C_2$. In particular
the original curves have no nonconstant elliptic
map or map to an abelian variety of two-rank zero.
\item Every geometric isogeny between $J_{C_0}$ and the
Jacobian of the quadratic Artin--Schreier twist of
the A.C2.13 split-branch curve has even degree.
\item The geometric Raynaud theta divisor has no
nonzero $\ell$-primary point for
\[
 \ell\in\{3,5,7,11,13,19,23,29,31,37,43,53,59,61,67\}.
                                                        \tag{C2.72}
\]
It has no point of order supported on any of
\[
 \{3,5\},\quad\{3,7\},\quad\{5,7\},\quad
 \{3,11\},\quad\{7,11\}.                              \tag{C2.73}
\]
Arbitrarily high powers of the listed primes are
included in these exclusions.
\item Its full geometric theta-$73$ locus consists of
four points on one cyclic subgroup. Relative to a
generator on theta, the packet is
\[
                    \{1,27,46,72\}=\{\pm1,\pm27\}.
                                                        \tag{C2.74}
\]
Only four of the eight primitive doubling orbits
meet theta. Every cyclic $73$-power cover therefore has
positive primitive two-rank; the degree-$73$ cover belonging
to this packet has primitive two-rank $36$ and
total two-rank $38$.
\end{enumerate}
These are exact limitations of specified
constructions. A1 for these two curves is not
decided by this claim.
\end{claim}
\begin{proof}
\proofparagraph{Counts and geometric simplicity.}
Direct Artin--Schreier enumeration on the smooth
normalization gives
\[
 \#C_0(\mathbf F_4)=7,\qquad
 \#C_0(\mathbf F_{16})=15,\qquad
 \#C_0(\mathbf F_{256})=255.
\]
At a simple rational pole there is one point;
elsewhere there are two precisely when the
absolute trace of the right side is zero.
The first two counts give
\[
           P(X)=X^4+2X^3+X^2+8X+16                     \tag{C2.75}
\]
for fourth-power Frobenius. The third count checks
its fourth-power trace. We use the exact
curve-to-Jacobian trace formula in
\cite[III, Theorem~11.1 and Corollary~11.4]{A_Milne2008},
together with the abelian-variety count and Weil
absolute values in its II, Theorem~1.1.
The three simple Artin--Schreier poles give
ordinariness.

Modulo three, $P$ is $\Phi_{10}$, irreducible
because $3$ has order four modulo $10$.
The squarefree nontrivial-root-ratio polynomial,
computed from $\operatorname{Res}_X(P(X),P(ZX))$
after deleting $(Z-1)^4$, is
\[
 \begin{split}
 R(Z)={}&64Z^8+80Z^7+108Z^6+7Z^5+26Z^4\\
        &\hspace{12mm}+7Z^3+108Z^2+80Z+64.
 \end{split}                                           \tag{C2.76}
\]
Its degree is at most eight for a structural
reason: write the four Weil roots as
$a,4/a,b,4/b$. There are four ratios within
reciprocal pairs and four cross ratios, the
latter each repeated twice. A root of unity
among these ratios would give a cyclotomic
factor of degree at most eight. The full list
of orders to check is
\[
 1,2,3,4,5,6,7,8,9,10,12,14,15,16,18,20,24,30.
\]
This list is exhaustive: only primes $2,3,5,7$
can divide such an order, with exponents at
most $4,2,1,1$, and the remaining finite
products give exactly the list. Exact
polynomial gcds with all eighteen cyclotomic
polynomials equal one.

Thus distinct roots of $P$ never become equal
after taking a positive power. Every $\pi^m$
has minimal polynomial of degree four, equal
to its irreducible characteristic polynomial.
A proper geometric abelian subvariety would
descend to some finite extension and factor
that polynomial. This proves geometric
simplicity.

\proofparagraph{The entire rational endomorphism algebra.}
Every geometric endomorphism commutes with
$\pi^m$ for some $m$. The tensor injection
\[
 \operatorname{End}(J)\otimes\mathbf Z_\ell
          \hookrightarrow\operatorname{End}(T_\ell J)
 \qquad(\ell\ne2)
\]
is \cite[I, Theorem~10.15, printed p.~49]{A_Milne2008}.
The centralizer of $\pi^m$ on the
four-dimensional rational Tate module is
$\mathbf Q_\ell[\pi^m]=\mathbf Q_\ell[\pi]$:
the operator has four distinct eigenvalues
and minimal polynomial of degree four.
All geometric endomorphisms therefore lie in
the same four-dimensional centralizer.
Their rational algebra contains the
degree-four field $\mathbf Q(\pi)$, so
equality follows. No assertion about a
maximal integral order is needed.

Put $s=\pi+4/\pi$. Equation (C2.75) gives
$s^2+2s-7=0$, and solving
$\pi^2-s\pi+4=0$ gives (C2.71).
Both real conjugates of $-7+4\sqrt2$ are
negative. Their product is $17$. If the CM
field were normal, it would contain both
square roots and hence $\sqrt{17}$.
Together with $\sqrt2$, this would generate
a totally real degree-four subfield, forcing
the whole field to be real, a contradiction.

An elliptic map would give an elliptic
Jacobian quotient. A nonconstant map to a
two-rank-zero abelian variety would give a
nonzero abelian image of the simple ordinary
Jacobian, hence an ordinary image of dimension
two. Both are impossible. Also the normalized
coefficients $(1,9,14)$ over
$\mathbf F_{16}=\mathbf F_2[w]/(w^4+w+1)$
are pairwise distinct. The six branch
permutations act by their coefficient
permutations, and an Artin--Schreier correction
cannot alter a simple-pole coefficient.
A.C2.22's classification therefore leaves
only the hyperelliptic involution.

\proofparagraph{Why a specific odd-isogeny transfer fails.}
The quadratic twist of the A.C2.13 curve
has the same polynomial (C2.75), as checked
independently by its counts $7,15,255$.
All three of its ramification points are
rational over $\mathbf F_4$, so Frobenius
fixes all three nonzero geometric two-torsion
classes. For $C_0$, Frobenius exchanges the
two finite ramification points and fixes
infinity. Its action on the three nonzero
classes is a transposition. Since $P(1)=28$,
the two rational two-primary groups are
respectively $C_2\times C_2$ and $C_4$.

For every $m$, the equality
$\mathbf Q(\pi^m)=\mathbf Q(\pi)$ gives
$R_m\in\mathbf Q[X]$ such that
$R_m(X^m)=X$ modulo $P(X)$.
The same polynomial works for both Jacobians.
A geometric homomorphism defined over
$\mathbf F_{4^m}$ intertwines their $m$-th
Frobenius powers; applying $R_m$ shows that
it intertwines their original fourth-power
Frobenius as well. An odd-degree isogeny
would give a Frobenius-equivariant isomorphism
of geometric two-torsion, contradicting the
two different actions. This concerns the
specified comparison pair, not all possible
isogeny correspondences or further covers.

\proofparagraph{Uniform torsion exponents from exact operator identities.}
If $P(X)\mid X^m-1$ modulo an odd prime
$\ell$, then $\pi^m$ fixes $J[\ell]$.
A.C2.12 makes every $\ell$-primary theta
point fixed by $\pi^m$. Suppose
\[
 (X^m-1)B(X)\equiv\ell^e
                  \pmod{P(X),\ell^{e+1}}.              \tag{C2.77}
\]
Because $P$ is monic, this yields an integral
endomorphism identity
\[
 (\pi^m-1)B(\pi)=\ell^e(1+\ell R_0(\pi)).
\]
On geometric $\ell$-primary torsion the
last factor is invertible, by its geometric
series at each finite torsion level. A
$\pi^m$-fixed point is therefore killed by
$\ell^e$. This does not presume that an
integral Frobenius lattice is cyclic.

Here are the exact certificates. The final
column lists the coefficients of $B$ in
ascending order; each row is checked directly
in the finite quotient ring in (C2.77).
\begin{center}
\scriptsize
\begin{tabular}{r r r l}
$\ell$&$m$&$\ell^e$&$B$ coefficients\\\hline
3&10&3&$(4,1,6,0)$\\
5&52&5&$(18,9,6,19)$\\
7&12&7&$(17,46,7,48)$\\
11&610&11&$(76,7,93,72)$\\
13&42&13&$(81,150,123,166)$\\
17&1224&289&$(1121,4202,390,3232)$\\
19&180&19&$(341,50,82,278)$\\
23&132&23&$(90,363,491,378)$\\
29&11788&29&$(425,801,273,143)$\\
31&120&31&$(379,567,530,11)$\\
37&24660&37&$(1116,23,450,215)$\\
41&140&1681&$(66658,64498,20991,27342)$\\
43&924&43&$(1238,65,550,591)$\\
47&23&2209&$(68200,85019,58876,86032)$\\
53&702&53&$(2277,11,505,1824)$\\
59&1740&59&$(851,1200,2335,1832)$\\
61&111660&61&$(3236,2301,2389,134)$\\
67&2244&67&$(201,915,1051,1020)$\\
71&1260&5041&$(157405,262215,277523,203927)$\\
73&1332&73&$(4521,1064,2697,1508)$
\end{tabular}
\end{center}
For every prime in (C2.72), the bound is
$\ell$ itself and the complete A.C2.18
Hankel test has residual maximal-minor gcd
one. This proves absence at every power.
The bounds at $17^2,41^2,47^2,71^2$ are
not asserted to be attained or excluded.

For a finite set $S$, first apply A.C2.14
to a common Frobenius power fixing every
$J[\ell]$, $\ell\in S$. It fixes the whole
$S$-supported theta point, hence every
primary component. Apply (C2.77) to these
components at that same power. The exact
component certificates give
\[
\begin{array}{c|r|r}
S&m&\text{order divides }N\\\hline
\{3,5\}&260&75\\
\{3,7\}&60&63\\
\{5,7\}&156&35\\
\{3,11\}&610&33\\
\{7,11\}&3660&77
\end{array}
\]
The complete geometric theta-$N$ test is
empty in each row. Thus (C2.73) includes
all powers, not only squarefree conductors.
The origin is outside the ordinary Raynaud
theta divisor. The exact ten component
polynomials and their remainders are in
\path{code/A_F4_remaining_mixed_packets.json}.

\proofparagraph{The complete theta-seventy-three packet.}
Use the binary encoding in A.C2.22, with
$w=2$ and $\alpha=w^2+w=6$ in $\mathbf F_{16}$.
The complete Hankel gcd is
\[
             T^2+6T+9=(T+8)(T+14).
\]
At these two coordinates the
Artin--Schreier constants are $9$ and $15$,
both of absolute trace one over $\mathbf F_{16}$.
Thus the four points lie over $\mathbf F_{256}$,
with sixteenth-power Frobenius exchanging
the two points above each coordinate.
The full norm witnesses verify order $73$;
alternatively A.C2.18's complete kernel
criterion gives annihilation by $73$, and
the non-Weierstrass condition makes the
class nonzero.

Equation (C2.75) gives
\[
                 \#J(\mathbf F_{256})=2^7\cdot7\cdot73.
\]
Its $73$-primary group is therefore cyclic
of order $73$, and contains all four theta
points. On normalized coordinates, actual
fourth-power Frobenius is
\[
                 (x,y)\longmapsto(x^4+1,y^4+\alpha).
\]
Its square is the hyperelliptic involution
on the four points. The theta origin class
$B_\theta$ is rational and satisfies
$2B_\theta=K_C$, so that involution acts by
negation on theta classes. Hence, if
$\pi\eta=k\eta$ for a theta generator,
$k^2=-1$ modulo $73$, giving $k=\pm27$.
The two coordinates are exchanged by
$\pi$, so the four points are exactly
$\pm\eta,\pm27\eta$.

The order of $2$ modulo $73$ is $9$.
The four multipliers in (C2.74) occupy
four distinct doubling orbits; four of
the eight remain. A.C2.5 gives primitive
two-rank $4\cdot9=36$, plus the base rank
two. Any other cyclic $73$ subgroup has
no theta point, so its primitive two-rank
is $72$. The exponent bound in the
$\ell=73$ row puts every $73$-primary
theta point in $J[73]$. Thus a cyclic
$73^a$-cover for $a\ge2$ has no primitive
theta character at all and has primitive
two-rank $\varphi(73^a)>0$.
This is a complete
$73$-primary statement, not merely a
rational-point sample. The separate Cantor
and scalar-Frobenius calculation reproduces
the same packet exactly.

\proofparagraph{Certificates and precise remaining scope.}
The author structural certificate is
\path{code/A_F4_remaining_structure.py};
the independent implementation
\path{code/A_F4_remaining_structure_degeneration.py}
uses its own finite-field and quotient-polynomial
arithmetic and checks the three counts,
all eighteen cyclotomic exclusions, and
all thirty primary and mixed identities.
The complete review is
\path{audit/A_F4_remaining_structure_degeneration_review.md}.
The complete conductor-$73$ gcd and norm
records are in
\path{code/A_theta_F4_remaining_73_89.json}
and
\path{code/A_F4_remaining_primitive_packets.py},
with the latter's adjacent JSON.
Root independently read the proof, reopened
the Milne passages, and reran the structural
and complete-packet certificates.

None of these assertions supplies colliding
sections. A.C2.25 below uses a different
four-prime conductor to settle A1 for both
curves, while every obstruction proved here
retains its stated scope.
\end{proof}
